\scanpage{241}
\[
\nu_p(n!)=\frac1{p-1}(n-\mathrm{chiff}_p n)
\]
\[
\nu_p(p^{(n)})=\nu_p\left(\frac{p^n}{n!}\right)=\frac1{p-1}\bigl((p-2)n+\uline{\mathrm{chiff}_p n}\bigr)\quad(=\mathrm{chiff}_p(n)\text{ si }p=2)
\]
\deleted{\unsure{Si on}}
\[
\begin{aligned}
\gamma_{p,i}&=\mathop{\mathrm{Inf}}_{\uline{n\geq i}}\nu_p(p^{(n)})=\mathop{\mathrm{Inf}}_{\uline{n\geq i}}\frac1{p-1}\underbrace{((p-2)n+\mathrm{chiff}_p(n))}_{\text{cette fonction en }n\text{ n'est ni croissante ni décroissante}\ldots}\\
&=\begin{cases}1&\text{si }p=2\\>\dfrac{p-2}{p-1}i\underset{i\to+\infty}{\longrightarrow}+\infty&\text{si }p\ne2\end{cases}
\end{aligned}
\]
\[
i=i_0+i_1p+\cdots+i_rp^r
\]
\deleted{\unsure{disons}}
\margin{\deleted{$i_0$}\quad\deleted{$\unsure{i}$}}
\[
\underbrace{\underbrace{(p-1)}_{i_0}+\cdots+\underbrace{(p-1)p^{s-1}}_{i_{s-1}}}_{\unsure{p^{s+1}-1|p^s-1}}+\underset{\unsure{<p-1}}{\underbrace{i_s}}p^{\unsure{s}}+\cdots+i_rp^r
\]
\[
i+1=\unsure{i_s+1|i_s}\,p^{\unsure{s}}+\cdots+i_rp^r
\]
\[
(p-2)i+(p-1)\unsure{s}\deleted{\unsure{0}}+\deleted{\sum_{\alpha\geq s}i_\alpha}\leq(p-2)(i+1)+\deleted{\sum_{\alpha\geq s}\unsure{i'_\alpha|i_\alpha}}
\]
\[
[\nu_p(p^{(i)})\leq\nu_p(p^{(i+1)})]\Longleftrightarrow[(p-1)s\leq(p-2)]\Longleftrightarrow[s=0]\Longleftrightarrow[i_0\ne p-1]
\]
\[
\deleted{\boxed{i_0\leq p-2}}\Longleftrightarrow i+1\not\equiv0\ (p)
\]
Exemple \unsure{où|si} $\gamma_{p,i}=\gamma_{p,i+1}$ : il \uline{suffit} qu'on ait \deleted{$\unsure{\nu_p(p)}$}
\[
\nu_p(p^{(i+1)})\leq\nu_p(p^{(i)})
\]
i.e. que \deleted{\unsure{dans le cas}} \deleted{\unsure{pour le cas}} $i+1\equiv0\ (p)$,
i.e. $i$ de la forme $p\ell-1$
\margin{240}
